\section{Analytical Properties of Leontovich Graphs}
\label{sec:leontovich}

\subsection{An eigenvalue condition that forbids crossovers}

\begin{theorem}[Single Positive Eigenvalue Reduction]
  \label{thm:spectral}
  Let $H$ be a rooted bipartite graph whose automorphism orbits respect
  the bipartition (i.e., no orbit contains vertices from both sides).
  If the quotient matrix of $H$ has exactly one positive eigenvalue,
  then the odd-$n$ near-path comparison has constant sign: for each fixed
  depth~$d$, the sign of $\Hom(P_n,H)-\Hom(E_n^{(d)},H)$ is determined on
  all odd~$n$ by any single verified odd base case.
\end{theorem}

\begin{proof}
  Since the orbits respect the bipartition, the quotient matrix has
  block form $\left[\begin{smallmatrix}0&B\\B^T&0\end{smallmatrix}\right]$
  after ordering the cells by side, so its nonzero spectrum is symmetric
  about zero.
  Because the base vector $\mathbf{1}$ and the initial degree vector
  are constant on automorphism orbits, the entire dynamic programming
  sequence $\{A^k \mathbf{1}\}$ remains in the column space of the
  quotient matrix, so the full adjacency spectrum is unnecessary.
  If there is exactly one positive eigenvalue $\lambda_1$, then every
  other nonzero eigenvalue is $-\lambda_1$, and the active spectrum is
  contained in $\{\lambda_1, 0, -\lambda_1\}$. By the quotient-eigenbasis
  expansion and Lemma~\ref{lem:active_recurrence},
  \[
    \Hom(G_n, H) = c\,\lambda_1^n + d\,(-\lambda_1)^n.
  \]
  For odd $n$, this reduces to $(c - d)\,\lambda_1^n$, and the kernel
  component vanishes once the exponent is at least $1$. There are no
  secondary exponential terms, so the ratio
  $\Hom(E_n, H) / \Hom(P_n, H)$ is constant on odd $n$. Thus the
  eigenvalue condition reduces any odd near-path crossover question to
  the sign of one verified odd base case in the family being studied.
\end{proof}

\begin{corollary}
  \label{cor:2_and_3_orbit_trees}
  All 2-orbit trees (stars $S_m$) and 3-orbit trees $T(x,y)$ have
  exactly one positive quotient eigenvalue. For these families, the
  eigenvalue condition reduces odd near-path crossover detection to a
  single base-case sign check.
\end{corollary}

\begin{corollary}[Double-cover transfer for small sweeps]
  \label{cor:loopy_transfer}
  If no connected simple graph on at most $2m$ vertices is Leontovich,
  then no connected graph on at most $m$ vertices, loops allowed, is
  Leontovich. In particular, if a connected looped target $H$ is
  non-bipartite, then its connected double cover $H \times K_2$ converts
  any unconditional simple-graph nonexistence bound at order~$\le 2m$ into
  the corresponding looped bound at order~$\le m$. Our bounded near-path
  sweeps therefore transfer only within their tested ranges: the
  simple-graph sweep up to $11$ vertices gives looped-target evidence on at
  most $5$ vertices for $n \le 200$, $d \le 20$, and the bipartite sweep
  up to $14$ vertices gives the analogous bounded evidence up to
  $7$ looped vertices.
\end{corollary}

\subsection{Four-orbit trees}

For 4-orbit trees $T(x,y,z)$, the characteristic polynomial of the
quotient matrix is
\begin{equation}
  \label{eq:charpoly}
  \lambda^4 - (x+y+z)\,\lambda^2 + xz = 0.
\end{equation}
Setting $s = x + y + z$, the discriminant of the quadratic in
$\lambda^2$ is
\[
  s^2 - 4xz = (x - z)^2 + y^2 + 2y(x + z).
\]
Since $x, y, z \ge 1$, this expression is always strictly positive
($\ge 5$). By Descartes' rule of signs, both roots for $\lambda^2$ are
positive, so \emph{every} spherically symmetric 4-orbit tree possesses
exactly 2 positive eigenvalues.

\subsection{The crossover threshold for \texorpdfstring{$T(7,1,9)$}{T(7,1,9)}}

This result was known privately to JD Nir and David Galvin prior to this work;
we record here the first fully written proof.

\begin{theorem}[Permanent Crossover]
  \label{thm:threshold}
  Let $H = T(7,1,9)$. Then $\Hom(E_n, H) < \Hom(P_n, H)$ for all
  odd $n \ge 13$, and direct computation gives $\Delta(11) < 0$ and
  $\Delta(13) > 0$.
\end{theorem}

\begin{proof}
  The characteristic polynomial of the quotient matrix
  \[
    M = \begin{pmatrix} 0&7&0&0 \\ 1&0&1&0 \\
    0&1&0&9 \\ 0&0&1&0 \end{pmatrix}
  \]
  is $t^4 - 17t^2 + 63 = 0$. Setting $u = t^2$:
  \[
    u = \frac{17 \pm \sqrt{37}}{2},
    \qquad
    \lambda_1 = \sqrt{\frac{17 + \sqrt{37}}{2}} \approx 3.3973,
    \qquad
    \lambda_2 = \sqrt{\frac{17 - \sqrt{37}}{2}} \approx 2.3364.
  \]
  For odd $n$, the quotient-eigenbasis expansion yields
  \[
    \Delta_{\mathrm{odd}}(n) \;=\; \Hom(P_n, H) - \Hom(E_n, H)
    \;=\; A\,\lambda_1^n - |B|\,\lambda_2^n,
  \]
  for constants $A > 0$ and $B \neq 0$. Setting
  $r = \lambda_2 / \lambda_1 = \sqrt{(17 - \sqrt{37})/(17 + \sqrt{37})}
  \approx 0.6877$:
  \[
    \Delta_{\mathrm{odd}}(n) = \lambda_1^n \bigl(A - |B|\,r^n\bigr).
  \]
  Since $r^n$ is strictly monotonically decreasing, $A - |B|\,r^n$ can
  cross zero at most once. Evaluating directly:
  \begin{align*}
    \Delta(11) &= 13{,}993{,}266 - 14{,}000{,}700 = -7{,}434 < 0, \\
    \Delta(13) &= 161{,}209{,}734 - 161{,}161{,}476 = +48{,}258 > 0.
  \end{align*}
  Thus $n = 13$ is the unique, permanent threshold.
\end{proof}

\begin{remark}[Parity heuristic, not theorem]
  \label{rem:parity}
  For even $n$, $P_n$ has a perfectly balanced bipartition $(n/2, n/2)$,
  minimizing its dominant projection. For even depths $d$ (e.g., $d=2$)
  where $E_n^{(d)}$ remains balanced, the path's linear geometry yields
  a smaller projection onto the dominant eigenvector. For odd depths
  (e.g., $d=1$), $E_n^{(d)}$ is forced into an unbalanced
  $(n/2+1, n/2-1)$ bipartition, giving it a strictly larger leading
  coefficient. This mechanism explains why even-$n$ crossovers were absent
  from many small sweeps, but Section~\ref{sec:doublecover} shows that it
  does not yield a universal obstruction. For odd $n$, both $P_n$ and $E_n^{(d)}$ share the same
  unbalanced bipartition $(\lceil n/2 \rceil, \lfloor n/2 \rfloor)$,
  enabling the secondary eigenvector to interfere and drive the crossover.
\end{remark}

\subsection{Minimum size in the symmetric family}

Exhaustive grid search over $T(x,y,z)$ with
$|V| = 1 + x + xy + xyz \le 100$ confirms that $T(7,1,9)$ with 78
vertices is the smallest Leontovich graph found in the $T(x,y,z)$
family within the tested range. A generalized search over all spherically symmetric trees
$T(d_1, \ldots, d_k)$ with $|V| \le 77$ and depth $k \le 6$ tested
5,732 trees with exact integer arithmetic
(\texttt{scripts/verify\_leontovich.py}, $n \le 200$, $d \le 20$) and
confirmed zero Leontovich instances (Table~\ref{tab:symmetric}).

\begin{table}[H]
\centering
\caption{Generalized symmetric tree search ($|V| \le 77$, exact integer
arithmetic, $n \le 200$, $d \le 20$).}
\label{tab:symmetric}
\begin{tabular}{@{}cccc@{}}
  \toprule
  Depth $k$ & Orbits & Trees tested & Leontovich found \\
  \midrule
  1 & 2 &    76 & \textbf{0} \\
  2 & 3 &   268 & \textbf{0} \\
  3 & 4 &   592 & \textbf{0} \\
  4 & 5 & 1,036 & \textbf{0} \\
  5 & 6 & 1,573 & \textbf{0} \\
  6 & 7 & 2,187 & \textbf{0} \\
  \midrule
  \multicolumn{2}{c}{Total} & 5,732 & \textbf{0} \\
  \bottomrule
\end{tabular}
\end{table}

\subsection{The \texorpdfstring{$n = 5$}{n=5} impossibility}

\begin{proposition}[$n = 5$ threshold is algebraically impossible]
  \label{prop:n5}
  For all $T(x,y,z)$ with $x, y, z \ge 1$,
  \[
    \Hom(P_5, T) - \Hom(E_5, T) = -xy\bigl((x - z - 1)^2 + y^2 z\bigr) < 0.
  \]
  In particular, $E_5$ always has \emph{more} homomorphisms than $P_5$
  into any $T(x,y,z)$, so no Leontovich crossover can occur at $n = 5$.
\end{proposition}

\begin{proof}
  Let $M$ denote the $4 \times 4$ quotient matrix of $T(x,y,z)$
  over its four automorphism orbits (root, children, grandchildren,
  leaves), with orbit-size vector
  $\mathbf{s} = (1,\, x,\, xy,\, xyz)^T$:
  \[
    M = \begin{pmatrix} 0&x&0&0 \\ 1&0&y&0 \\
    0&1&0&z \\ 0&0&1&0 \end{pmatrix}.
  \]
  For any path $P_n$, the homomorphism count decomposes as
  $\Hom(P_n, H) = \mathbf{s}^T M^{n-1} \mathbf{1}$,
  where $\mathbf{1} = (1,1,1,1)^T$. Iterating:
  \begin{align*}
    M\,\mathbf{1} &= (x,\; y+1,\; z+1,\; 1)^T, \\
    M^2\mathbf{1} &= (xy+x,\; x+yz+y,\; y+z+1,\; z+1)^T.
  \end{align*}
  The tree $E_5$ has vertices $\{1,2,3,4,5\}$ with edges
  $(1,2),\,(2,3),\,(3,4),\,(2,5)$, so vertex~2 has degree~3. Rooting
  at vertex~2 and summing over orbit assignments:
  \[
    \Hom(E_5, H) = \sum_{b=0}^{3}
      s_b \cdot (M\mathbf{1})_b^2 \cdot (M^2\mathbf{1})_b,
  \]
  since the three branches from vertex~2 contribute
  one factor of $(M\mathbf{1})_b$ each for the two pendant
  leaves (vertices~1 and~5) and one factor of
  $(M^2\mathbf{1})_b$ for the length-2 branch (vertices 3--4).

  Expanding $\mathbf{s}^T M^4 \mathbf{1}$ and collecting terms:
  \begin{align*}
    \Hom(P_5, T) - \Hom(E_5, T)
    &= -xy\bigl(x^2 - 2x(z+1) + (z+1)^2 + y^2 z\bigr) \\
    &= -xy\bigl((x - z - 1)^2 + y^2 z\bigr).
  \end{align*}
  Since $x, y, z \ge 1$, the squared term
  $(x - z - 1)^2 \ge 0$ and $y^2 z \ge 1$, so the
  parenthesised sum is strictly positive. Multiplied by $-xy < 0$,
  the difference is strictly negative.
\end{proof}

\begin{remark}[The remaining $n=6$ cases]
  \label{rem:n6_residue}
  The six trees on six vertices split into four proved cases and two
  checked cases. The path is the reference case; the star $K_{1,5}$
  follows from Sidorenko's theorem; the balanced near-path $\Sp(2,2,1)$
  is controlled by Theorem~\ref{thm:spectral}; and the double star
  $S_{2,2}$ is covered by the Cauchy--Schwarz bound above.
  The remaining two six-vertex trees, $\Sp(3,1,1)$ and
  $\Sp(2,1,1,1)$, are supported only by exact computation. Consequently,
  the four proved cases carry label (P), while these two cases carry label
  (E). A proof covering all six cases would require resolving this gap; the
  near-path sufficiency conjecture also remains open.
\end{remark}

\subsection{Minimum threshold landscape}

Sweeping the generator over all $T(x,y,z)$ with $|V| \le 500$ and
selecting the smallest graph achieving each odd-$n$ threshold yields
Table~\ref{tab:thresholds}. The vertex count bottoms out at $n = 13$
with $T(7,1,9)$, then \emph{increases} again---the threshold at $n = 15$
requires a larger graph than $n = 13$.

\begin{table}[H]
\centering
\caption{Smallest Leontovich graph by crossover threshold within the tested
range.}
\label{tab:thresholds}
\begin{tabular}{@{}llcc@{}}
  \toprule
  Family & Parameters & $|V|$ & Threshold \\
  \midrule
  $T(x,y,z)$ & $T(16, 1, 19)$ & 337 & $n = 7$ \\
              & $T(9, 1, 11)$  & 118 & $n = 9$ \\
              & $T(8, 1, 10)$  &  97 & $n = 11$ \\
              & $T(7, 1, 9)$   &  78 & $n = 13$ \\
              & $T(8, 1, 12)$  & 113 & $n = 15$ \\
  \bottomrule
\end{tabular}
\end{table}

\begin{remark}
  Galvin et al.~\cite{galvin2025} used $T(18,3,32)$ with $1{,}801$ vertices to prove
  the $n = 7$ crossover (their Theorem~1.3). Our landscape sweep shows
  that $T(16,1,19)$ achieves the identical crossover with only
  $337$~vertices---a $5.3\times$ reduction.  This is the smallest
  $T(x,y,z)$ graph achieving a threshold $n \le 7$ within the tested range.
\end{remark}

% ════════════════════════════════════════════════════════════════════════
